\documentclass[12pt]{article}
\usepackage{amsmath, amssymb, amsthm}
\usepackage[margin=1in]{geometry}
\usepackage{hyperref}

\title{Problem 254: Sums of Digit Factorials}
\author{Project Euler Solution}
\date{}

\begin{document}
\maketitle

\section{Problem Statement}
Define:
\begin{align}
f(n) &= \text{sum of factorials of digits of } n \\
sf(n) &= \text{digit sum of } f(n) \\
g(i) &= \min\{n : sf(n) = i\} \\
sg(i) &= \text{digit sum of } g(i)
\end{align}

Find $\displaystyle\sum_{i=1}^{150} sg(i)$.

\section{Key Observations}

\begin{enumerate}
    \item $f(n)$ depends only on the multiset of digits of $n$.
    \item The smallest $n$ with a given digit multiset has digits in non-decreasing order.
    \item For large $i$, $g(i)$ consists of many 9's plus a small adjustment set of digits.
    \item $sf(n) = s(f(n))$ where $s$ is the digit sum function.
\end{enumerate}

\section{Algorithm}

For each target $i$ from 1 to 150:
\begin{enumerate}
    \item Enumerate ``base patterns'' (multisets of digits from 1--9 with at most a few non-9 digits).
    \item For each pattern, determine the number of 9's needed so that $s(f(n)) = i$.
    \item The digit sum of $g(i)$ is $sg(i) = 9 \cdot c_9 + \sum_{d} d \cdot c_d$ for the optimal pattern.
\end{enumerate}

The number of 9's $c_9$ is chosen so that $f(n) = c_9 \cdot 362880 + R$ (where $R$ is the sum of factorials of the non-9 digits) has digit sum equal to $i$.

\section{Result}
\[
\boxed{8184523820510}
\]

\section{Answer}

\[
\boxed{8184523820510}
\]

\end{document}
